\section*{Executive Summary}
\addcontentsline{toc}{section}{Executive Summary}

For over three decades, management of western monarch butterfly overwintering habitat in California has been guided by the Disruptive Wind Hypothesis, which holds that wind speeds at or above 2 m/s (approximately 5 mph) physically dislodge butterflies from their roosting clusters or force them to abandon sites. This threshold has been widely adopted in management guidelines, restoration planning, and land use decisions affecting overwintering groves, directing substantial conservation investment toward wind protection at overwintering sites.

Despite its influence on management practice, the wind disruption threshold had never been directly tested through empirical observation of butterfly behavioral responses to wind. This study, conducted at Vandenberg Space Force Base during the 2023--2024 overwintering season, provides the first such test. Using time-lapse cameras and wind sensors deployed at monarch cluster roost sites, we collected 1,894 paired observations over 78 monitoring days, enabling direct analysis of the relationship between wind conditions and butterfly clustering behavior.

The results do not support the wind disruption hypothesis. Every observation period in our day-to-day analysis experienced wind speeds above the 2 m/s threshold---often by a wide margin---yet monarch clusters persisted throughout the monitoring period. Statistical analysis found no meaningful relationship between wind speed and changes in cluster size at either short-term (30-minute) or daily timescales. When butterflies remained in shade, wind had no detectable effect on cluster dynamics across the entire observed range (0--12.4 m/s). Statistical power analysis confirmed that the study had sufficient power (87.5\% for moderate effects) to detect biologically meaningful wind effects if they existed.

Instead of wind disruption, analyses consistently identified interactions between wind and sunlight as significant predictors of cluster dynamics, suggesting that thermal regulation---not wind avoidance---may be the primary driver of butterfly movement at roost sites. These findings align with recent work questioning other components of the original microclimate hypothesis and pointing toward canopy structure and light patterns as key habitat features.

\subsection*{Management Implications}

These findings carry several implications for the management of overwintering habitat. The 2 m/s wind speed threshold does not appear to represent a meaningful biological boundary for overwintering monarchs and should not be used as a sole criterion for habitat suitability assessment. Suitable overwintering habitat within existing groves may be more extensive than currently recognized, as areas previously excluded based on wind exposure may support clusters if they provide appropriate thermal and light conditions. Correspondingly, management requirements for wind buffer zones---sometimes extending to six tree heights (approximately 240 m) from cluster locations---may be more restrictive than necessary.

Existing trees and grove structure should be maintained as a precaution, but the rationale for their protection should be updated to reflect current evidence. Future management and restoration efforts should prioritize understanding canopy structure and the resulting light patterns, which appear more directly connected to monarch site selection than wind speed alone.
